% chapterintro
% ==============================
% 第14章 导言：14.0-chinese.md
本章的主要内容有：

\begin{itemize}
\item 将测试部署到CDash面板

\item CDash面板显示测试覆盖率

\item 使用AddressSanifier向CDash报告内存缺陷

\item 使用ThreadSaniiser向CDash报告数据争用

\end{itemize}

CDash是一个web服务，用于汇集CTest在测试运行期间、夜间测试期间或在持续集成中的测试结果。面板报告就是我们所说的\textbf{CDash时}，如下图所示：

\myGraphic{0.9}{images/preface/2.png}{}

本章中，我们将向CDash报告测试结果。将讨论报告测试覆盖率的策略，以及分别使用AddressSanifier和ThreadSanifier等工具，收集的内存缺陷和数据争用问题。

有两种方法向CDash报告结果：

\begin{enumerate}
\item 通过构建的测试目标

\item 使用CTest脚本

\end{enumerate}

在前两个示例中使用建立测试目标的方式，在后两个示例中使用CTest脚本。

\mySmallsection{CDash环境}

CDash的安装需要使用PHP和SSL的web服务器(Apache、NGINX或IIS)，并访问MySQL或PostgreSQL数据库服务器。详细讨论CDash web服务的设置超出了本书的范围，读者们可以参考官方文档：\texthref{https://public.kitware.com/Wiki/CDash:Installation}

Kitware提供了两个面板(https://my.cdash.org 和 \texthref{https://open.cdash.org} )，因此本章中的示例并不需要安装CDash。我们将在示例中参考已经提供的面板。

对于想要自己安装CDash的读者，我们建议使用MySQL作为后端，因为这是 \texthref{https://my.cdash.org} 和 \texthref{https://open.cdash.org} 的配置方式，而且社区也对这种搭配方式进行了测试。

\begin{myInfoBox}
\textbf{NOTE}：也可以使用Docker来安装CDash。官方镜像的请求在CDash的跟踪器上处于打开状态，网址是\texthref{https://github.com/Kitware/CDash/issues/562}
\end{myInfoBox}


% --- recipe 1: 14.1-chinese.md ---
\section{14.1 将测试部署到CDash}

\begin{myInfoBox}
\textbf{NOTE}：此示例代码可以在 \texthref{https://github.com/dev-cafe/cmake-cookbook/tree/v1.0/chapter-14/recipe-01} 中找到，其中包含一个C++示例。该示例在CMake 3.5版(或更高版本)中是有效的，并且已经在GNU/Linux、macOS和Windows上进行过测试。
\end{myInfoBox}

本示例中，我们将扩展第4章第1节的测试示例，并将测试结果部署到\texthref{https://my.cdash.org/index.php?project=cmake-cookbook} ，这是在Kitware为社区提供的公共面板( \texthref{https://my.cdash.org} )的基础上，为本书创建的专属面板。

\mySmallsection{准备工作}

我们将从重用第1节中的示例源代码，该测试将整数作为命令行参数进行求和。该示例由三个源文件组成:\texttt{main.cpp}、\texttt{sum\_integer.cpp}和\texttt{sum\_integers.hpp}。我们还将重用第4章(创建和运行测试)中的\texttt{test.cpp}文件，但这里将它重命名为\texttt{test\_short.cpp}。我们将使用\texttt{test\_long.cpp}扩展这个例子:

\begin{codebox}[cpp]
#include "sum_integers.hpp"

#include <numeric>
#include <vector>

int main() {
  
  // creates vector {1, 2, 3, ..., 999, 1000}
  std::vector integers(1000);
  std::iota(integers.begin(), integers.end(), 1);

    if (sum_integers(integers) == 500500) {
    return 0;
  } else {
    return 1;
  }
}
\end{codebox}

然后，将这些文件组织成以下文件树:

\begin{codebox}[shell]
.
├── CMakeLists.txt
├── CTestConfig.cmake
├── src
│    ├── CMakeLists.txt
│    ├── main.cpp
│    ├── sum_integers.cpp
│    └── sum_integers.hpp
└── tests
    ├── CMakeLists.txt
    ├── test_long.cpp
    └── test_short.cpp
\end{codebox}

\mySmallsection{具体实施}

现在，我们将演示如何配置、构建、测试。最后，将示例项目的测试结果提交到面板的过程：

\begin{enumerate}
\item 源目标在\texttt{src/CMakeLists.txt}中定义，如下：

\begin{codebox}[cmake]
# example library
add_library(sum_integers "")
target_sources(sum_integers
  PRIVATE
    sum_integers.cpp
  PUBLIC
    ${CMAKE_CURRENT_LIST_DIR}/sum_integers.hpp
  )
target_include_directories(sum_integers
  PUBLIC
    ${CMAKE_CURRENT_LIST_DIR}
  )
# main code
add_executable(sum_up main.cpp)
target_link_libraries(sum_up sum_integers)
\end{codebox}

\item \texttt{tests/CMakeLists.txt}中定义了测试：

\begin{codebox}[cmake]
add_executable(test_short test_short.cpp)
target_link_libraries(test_short sum_integers)
add_executable(test_long test_long.cpp)
target_link_libraries(test_long sum_integers)
add_test(
  NAME
    test_short
  COMMAND
    $<TARGET_FILE:test_short>
  )
add_test(
  NAME
    test_long
  COMMAND
    $<TARGET_FILE:test_long>
  )
\end{codebox}

\item 主\texttt{CMakeLists.txt}文件引用前面的两个文件，这个配置中的新元素是\texttt{include(CTest)}，这样就可以向CDash仪表板报告结果：

\begin{codebox}[cmake]
# set minimum cmake version
cmake_minimum_required(VERSION 3.5 FATAL_ERROR)
# project name and language
project(recipe-01 LANGUAGES CXX)
# require C++11
set(CMAKE_CXX_STANDARD 11)
set(CMAKE_CXX_EXTENSIONS OFF)
set(CMAKE_CXX_STANDARD_REQUIRED ON)
# process src/CMakeLists.txt
add_subdirectory(src)
enable_testing()
# allow to report to a cdash dashboard
include(CTest)
# process tests/CMakeLists.txt
add_subdirectory(tests)
\end{codebox}

\item 另外，我们创建文件\texttt{CTestConfig.cmake}与主\texttt{CMakeLists.txt}文件位于同一目录中。这个新文件包含以下几行：

\begin{codebox}[cmake]
set(CTEST_DROP_METHOD "http")
set(CTEST_DROP_SITE "my.cdash.org")
set(CTEST_DROP_LOCATION "/submit.php?project=cmake-cookbook")
set(CTEST_DROP_SITE_CDASH TRUE)
\end{codebox}

\item 我们现在已经准备好配置和构建项目：

\begin{codebox}[shell]
$ mkdir -p build
$ cd build
$ cmake ..
$ cmake --build .
\end{codebox}

\item 构建后，运行测试集，并向面板报告测试结果：

\begin{codebox}[shell]
$ ctest --dashboard Experimental
Site: larry
Build name: Linux-c++
Create new tag: 20180408-1449 - Experimental
Configure project
Each . represents 1024 bytes of output
. Size of output: 0K
Build project
Each symbol represents 1024 bytes of output.
'!' represents an error and '*' a warning.
. Size of output: 0K
0 Compiler errors
0 Compiler warnings
Test project /home/user/cmake-recipes/chapter-15/recipe-01/cxx-example/build
Start 1: test_short
1/2 Test #1: test_short ....................... Passed 0.00 sec
Start 2: test_long
2/2 Test #2: test_long ........................ Passed 0.00 sec
100% tests passed, 0 tests failed out of 2
Total Test time (real) = 0.01 sec
Performing coverage
Cannot find any coverage files. Ignoring Coverage request.
Submit files (using http)
Using HTTP submit method
Drop site:http://my.cdash.org/submit.php?project=cmake-cookbook
Uploaded: /home/user/cmake-recipes/chapter-14/recipe-01/cxx-example/build/Testing/20180408-1449/Build.xml
Uploaded: /home/user/cmake-recipes/chapter-14/recipe-01/cxx-example/build/Testing/20180408-1449/Configure.xml
Uploaded: /home/user/cmake-recipes/chapter-14/recipe-01/cxx-example/build/Testing/20180408-1449/Test.xml
Submission successful
\end{codebox}

\item 最后，可以在浏览器中看到测试结果(本例中，测试结果上报到 \texthref{https://my.cdash.org/index.php?project=cmake-cookbook} ):

\myGraphic{0.9}{images/chapter14/14-1.png}{}

\end{enumerate}

\mySmallsection{工作原理}

可以从更高级的角度展示工作流，CTest运行测试并在XML文件中记录结果。然后，将这些XML文件发送到CDash服务器，在那里可以浏览和分析它们。通过单击数字\texttt{2}，获得关于通过或失败测试的更多的细节信息(本例中，没有失败的测试)。如下图所示，详细记录了运行测试的机器的信息，以及时间信息。同样，单个测试的测试输出也可以在线浏览。

\myGraphic{0.9}{images/chapter14/14-2.png}{}

CTest支持三种不同的提交模式：

\begin{itemize}
\item 实验性构建

\item 夜间构建

\item 持续构建

\end{itemize}

我们使用了\texttt{ ctest --dashboard Experimental }(实验性构建提交)，因此，测试结果显示在实验模式之下。实验模式对于测试代码的当前状态、调试新的仪表板脚本、调试CDash服务器或项目非常有用。夜间构建模式，将把代码更新(或降级)到最接近最近夜间构建开始时的存储库，这些可以在\texttt{ CTestConfig.cmake }中设置。其为接收更新频繁的项目的所有夜间测试提供一个定义良好的参考。例如，夜间开始时间可以设置为世界时的"午夜"：

\begin{codebox}[cmake]
set(CTEST_NIGHTLY_START_TIME "00:00:00 UTC")
\end{codebox}

持续模式对于集成工作流非常有用，它将把代码更新到最新版本。

\begin{myInfoBox}
\textbf{TIPS}：构建、测试和提交到实验面板只需要一个命令—\texttt{ cmake --build . --target Experimental}
\end{myInfoBox}

\mySmallsection{更多信息}

这个示例中，我们直接从测试目标部署到CDash。我们将在本章后面的第3和第4部分中，使用专用的CTest脚本。

CDash不仅可以监视测试是否通过或失败，还可以看到测试时间。可以为测试计时进行配置：如果测试花费的时间超过分配的时间，它将被标记为失败。这对于基准测试非常有用，可以在重构代码时自动检测性能测试用例的性能情况。

有关CDash定义和配置设置的详细讨论，请参见官方CDash文档，网址为 \texthref{https://public.kitware.com/Wiki/CDash:Documentation}


% --- recipe 2: 14.2-chinese.md ---
\section{14.2 CDash显示测试覆盖率}

\begin{myInfoBox}
\textbf{NOTE}：此示例代码可以在 \texthref{https://github.com/dev-cafe/cmake-cookbook/tree/v1.0/chapter-14/recipe-02} 中找到，其中包含一个C++示例。该示例在CMake 3.5版(或更高版本)中是有效的，并且已经在GNU/Linux、macOS和Windows上进行过测试。
\end{myInfoBox}

本示例中，我们将测试覆盖率报告给CDash，面板上将能够逐行浏览测试覆盖率分析，以便识别未测试或未使用的代码。

\mySmallsection{准备工作}

我们将扩展前一节的源代码，在\texttt{src/sum\_integers.cpp}中做一个小的修改，添加一个函数\texttt{ sum\_integers\_unused }:

\begin{codebox}[cpp]
#include "sum_integers.hpp"

#include <vector>

int sum_integers(const std::vector integers) {
  auto sum = 0;
  
  for (auto i : integers) {
    sum += i;
  }
  
  return sum;
}

int sum_integers_unused(const std::vector integers) {
  auto sum = 0;
  
  for (auto i : integers) {
    sum += i;
  }
  
  return sum;
}
\end{codebox}

我们使用gcov(https://gcc.gnu.org/onlinedocs/gcc/Gcov.html )通过覆盖率分析检测这个未使用的代码。

\mySmallsection{具体实施}

通过以下步骤，我们将使用覆盖率分析，并将结果上传到面板：

\begin{enumerate}
\item 主\texttt{CMakeLists.txt}和\texttt{tests/CMakeLists.txt}文件与前一个示例相同。

\item 我们将扩展\texttt{src/CMakeLists.txt}，并提供一个选项来添加用于代码覆盖率的编译标志。此选项默认启用:

\begin{codebox}[cmake]
option(ENABLE_COVERAGE "Enable coverage" ON)
if(ENABLE_COVERAGE)
  if(CMAKE_CXX_COMPILER_ID MATCHES GNU)
    message(STATUS "Coverage analysis with gcov enabled")
    target_compile_options(sum_integers
      PUBLIC
        -fprofile-arcs -ftest-coverage -g
      )
    target_link_libraries(sum_integers
      PUBLIC
        gcov
      )
  else()
    message(WARNING "Coverage not supported for this compiler")
  endif()
endif()
\end{codebox}

\item 然后，配置、构建，并将结果上传CDash:

\begin{codebox}[shell]
$ mkdir -p build
$ cd build
$ cmake ..
$ cmake --build . --target Experimental
\end{codebox}

\item 最后一步，执行测试覆盖率分析:

\begin{codebox}[cmake]
  Performing coverage
  Processing coverage (each . represents one file):
  ...
  Accumulating results (each . represents one file):
  ...
    Covered LOC: 14
    Not covered LOC: 7
    Total LOC: 21
    Percentage Coverage: 66.67%
Submit files (using http)
  Using HTTP submit method
  Drop site:http://my.cdash.org/submit.php?project=cmake-cookbook
  Uploaded: /home/user/cmake-recipes/chapter-14/recipe-02/cxx-example/build/Testing/20180408-1530/Build.xml
  Uploaded: /home/user/cmake-recipes/chapter-14/recipe-02/cxx-example/build/Testing/20180408-1530/Configure.xml
  Uploaded: /home/user/cmake-recipes/chapter-14/recipe-02/cxx-example/build/Testing/20180408-1530/Coverage.xml
  Uploaded: /home/user/cmake-recipes/chapter-14/recipe-02/cxx-example/build/Testing/20180408-1530/CoverageLog-0.xml
  Uploaded: /home/user/cmake-recipes/chapter-14/recipe-02/cxx-example/build/Testing/20180408-1530/Test.xml
  Submission successful
\end{codebox}

\item 最后，可以在浏览器中验证测试结果(本例的测试结果报告在 \texthref{https://my.cdash.org/index.php?project=cmake-cookbook} ):

\myGraphic{0.9}{images/chapter14/14-3.png}{}

\end{enumerate}

\mySmallsection{工作原理}

测试覆盖率为66.67\%。为了得到更深入的了解，我们可以点击百分比，得到两个子目录的覆盖率分析:

\myGraphic{0.9}{images/chapter14/14-4.png}{}

通过浏览子目录链接，我们可以检查单个文件的测试覆盖率，甚至可以逐行浏览摘要(例如，\texttt{src/sum\_integs.cpp})：

\myGraphic{0.9}{images/chapter14/14-5.png}{}

运行测试时，绿线部分已经被覆盖，而红线部分则没有。通过这个方法，我们不仅可以标识未使用的/未测试的代码(使用\texttt{sum\_integers\_used}函数)，还可以查看每一行代码被遍历的频率。例如，代码行\texttt{sum += i}已经被访问了1005次(在\texttt{test\_short}期间访问了5次，在\texttt{test\_long}期间访问了1000次)。测试覆盖率分析是自动化测试不可或缺的功能，CDash为我们提供了一个界面，可以在浏览器中图形化地浏览分析结果。

\mySmallsection{更多信息}

为了更多的了解该特性，我们推荐读者阅读下面的博客文章，它更深入的讨论了CDash的覆盖特性：\texthref{https://blog.kitware.com/additional-coverage-features-in-cdash/}


% --- recipe 3: 14.3-chinese.md ---
\section{14.3 使用AddressSanifier向CDash报告内存缺陷}

\begin{myInfoBox}
\textbf{NOTE}：此示例代码可以在 \texthref{https://github.com/dev-cafe/cmake-cookbook/tree/v1.0/chapter-14/recipe-03} 中找到，其中包含一个C++示例和一个Fortran例子。该示例在CMake 3.5版(或更高版本)中是有效的，并且已经在GNU/Linux、macOS和Windows上进行过测试。
\end{myInfoBox}

AddressSanitizer(ASan)是可用于C++、C和Fortran的内存检测。它可以发现内存缺陷，比如：在空闲后使用、返回后使用、作用域后使用、缓冲区溢出、初始化顺序错误和内存泄漏(请参见 \texthref{https://github.com/google/sanitizers/wiki/AddressSanitizer} )。从3.1版本开始，AddressSanitizer是LLVM的一部分；从4.8版本开始，作为GCC的一部分。在这个示例中，我们将在代码中加入两个bug，正常的测试中可能无法检测到。为了检测这些bug，我们将使用AddressSanitizer工具，并将CTest与动态分析结合起来，从而将缺陷报告给CDash。

\mySmallsection{准备工作}

这个例子中，我们将使用两个源文件和两个测试集：

\begin{codebox}[shell]
.
├── CMakeLists.txt
├── CTestConfig.cmake
├── dashboard.cmake
├── src
│    ├── buggy.cpp
│    ├── buggy.hpp
│    └── CMakeLists.txt
└── tests
    ├── CMakeLists.txt
    ├── leaky.cpp
    └── use_after_free.cpp
\end{codebox}

\texttt{buggy.cpp}包含有两个bug：

\begin{codebox}[cpp]
#include "buggy.hpp"

#include <iostream>

int function_leaky() {
  double *my_array = new double[1000];
  // do some work ...
  // we forget to deallocate the array
  // delete[] my_array;
  return 0;
}

int function_use_after_free() {
  double *another_array = new double[1000];
  // do some work ...
  // deallocate it, good!
  delete[] another_array;
  // however, we accidentally use the array
  // after it has been deallocated
  std::cout << "not sure what we get: " << another_array[123] << std::endl;
  return 0;
}
\end{codebox}

这些函数在相应的头文件中声明(\texttt{buggy.hpp}）：

\begin{codebox}[cpp]
#pragma once
int function_leaky();
int function_use_after_free();
\end{codebox}

测试文件\texttt{leaky.cpp}中将会验证\texttt{function\_leaky}的返回值：

\begin{codebox}[cpp]
#include "buggy.hpp"
int main() {
  int return_code = function_leaky();
  return return_code;
}
\end{codebox}

相应地，\texttt{use\_after\_free.cpp}会检查\texttt{function\_use\_after\_free}的返回值:

\begin{codebox}[cpp]
#include "buggy.hpp"
int main() {
  int return_code = function_use_after_free();
  return return_code;
}
\end{codebox}

\mySmallsection{具体实施}

为了使用ASan，我们需要使用特定的标志来编译代码。然后，我们将运行测试并将它们提交到面板。

\begin{enumerate}
\item 生成bug库的工作将在\texttt{ src/CMakeLists.txt}中完成：

\begin{codebox}[cmake]
add_library(buggy "")
target_sources(buggy
  PRIVATE
    buggy.cpp
  PUBLIC
    ${CMAKE_CURRENT_LIST_DIR}/buggy.hpp
  )
target_include_directories(buggy
  PUBLIC
    ${CMAKE_CURRENT_LIST_DIR}
  )
\end{codebox}

\item 在文件\texttt{src/CMakeLists.txt}中，我们将添加一个选项用于使用ASan：

\begin{codebox}[cmake]
option(ENABLE_ASAN "Enable AddressSanitizer" OFF)
if(ENABLE_ASAN)
  if(CMAKE_CXX_COMPILER_ID MATCHES GNU)
    message(STATUS "AddressSanitizer enabled")
    target_compile_options(buggy
      PUBLIC
        -g -O1 -fsanitize=address -fno-omit-frame-pointer
      )
    target_link_libraries(buggy
      PUBLIC
        asan
      )
  else()
    message(WARNING "AddressSanitizer not supported for this compiler")
  endif()
endif()
\end{codebox}

\item 测试在\texttt{tests/CMakeLists.txt}中定义:

\begin{codebox}[cmake]
foreach(_test IN ITEMS leaky use_after_free)
  add_executable(${_test} ${_test}.cpp)
  target_link_libraries(${_test} buggy)
  add_test(
    NAME
        ${_test}
    COMMAND
        $<TARGET_FILE:${_test}>
    )
endforeach()
\end{codebox}

\item 主\texttt{CMakeLists.txt}与之前的示例基本相同：

\begin{codebox}[cmake]
# set minimum cmake version
cmake_minimum_required(VERSION 3.5 FATAL_ERROR)
# project name and language
project(recipe-03 LANGUAGES CXX)
# require C++11
set(CMAKE_CXX_STANDARD 11)
set(CMAKE_CXX_EXTENSIONS OFF)
set(CMAKE_CXX_STANDARD_REQUIRED ON)
# process src/CMakeLists.txt
add_subdirectory(src)
enable_testing()
# allow to report to a cdash dashboard
include(CTest)
# process tests/CMakeLists.txt
add_subdirectory(tests)
\end{codebox}

\item \texttt{ CTestConfig.cmake }也没有修改：

\begin{codebox}[cmake]
set(CTEST_DROP_METHOD "http")
set(CTEST_DROP_SITE "my.cdash.org")
set(CTEST_DROP_LOCATION "/submit.php?project=cmake-cookbook")
set(CTEST_DROP_SITE_CDASH TRUE)
\end{codebox}

\item 这个示例中，我们使用CTest脚本向CDash提交结果；为此，我们将创建一个文件\texttt{ dashboard.cmake}(与主\texttt{CMakeLists.txt}和`\texttt{ CTestConfig.cmake}位于同一个目录下)：

\begin{codebox}[cmake]
set(CTEST_PROJECT_NAME "example")
cmake_host_system_information(RESULT _site QUERY HOSTNAME)
set(CTEST_SITE ${_site})
set(CTEST_BUILD_NAME "${CMAKE_SYSTEM_NAME}-${CMAKE_HOST_SYSTEM_PROCESSOR}")
set(CTEST_SOURCE_DIRECTORY "${CTEST_SCRIPT_DIRECTORY}")
set(CTEST_BINARY_DIRECTORY "${CTEST_SCRIPT_DIRECTORY}/build")
include(ProcessorCount)
ProcessorCount(N)
if(NOT N EQUAL 0)
  set(CTEST_BUILD_FLAGS -j${N})
  set(ctest_test_args ${ctest_test_args} PARALLEL_LEVEL ${N})
endif()
ctest_start(Experimental)
ctest_configure(
  OPTIONS
    -DENABLE_ASAN:BOOL=ON
  )
ctest_build()
ctest_test()
set(CTEST_MEMORYCHECK_TYPE "AddressSanitizer")
ctest_memcheck()
ctest_submit()
\end{codebox}

\item 我们将执行\texttt{ dashboard.cmake}脚本。注意，我们使用\texttt{CTEST\_CMAKE\_GENERATOR}与生成器选项的方式：

\begin{codebox}[shell]
$ ctest -S dashboard.cmake -D 
CTEST_CMAKE_GENERATOR="Unix Makefiles"
Each . represents 1024 bytes of output
. Size of output: 0K
Each symbol represents 1024 bytes of output.
'!' represents an error and '*' a warning.
. Size of output: 1K
\end{codebox}

\item 结果将会出现在CDash网站上:

\myGraphic{0.9}{images/chapter14/14-6.png}{}

\end{enumerate}

\mySmallsection{具体实施}

这个示例中，成功地向仪表板的动态分析部分报告了内存错误。我们可以通过浏览缺陷详细信息，得到进一步的了解:

\myGraphic{0.9}{images/chapter14/14-7.png}{}

通过单击各个链接，可以浏览完整信息的输出。

注意，也可以在本地生成AddressSanitizer报告。这个例子中，我们需要设置\texttt{ENABLE\_ASAN}:

\begin{codebox}[shell]
$ mkdir -p build
$ cd build
$ cmake -DENABLE_ASAN=ON ..
$ cmake --build .
$ cmake --build . --target test

Start 1: leaky
1/2 Test #1: leaky ............................***Failed 0.07 sec
Start 2: use_after_free
2/2 Test #2: use_after_free ...................***Failed 0.04 sec
0% tests passed, 2 tests failed out of 2
\end{codebox}

运行\texttt{leaky}测试，直接产生以下结果:

\begin{codebox}[shell]
$ ./build/tests/leaky

=================================================================
==18536==ERROR: LeakSanitizer: detected memory leaks
Direct leak of 8000 byte(s) in 1 object(s) allocated from:
#0 0x7ff984da1669 in operator new[](unsigned long) /build/gcc/src/gcc/libsanitizer/asan/asan_new_delete.cc:82
#1 0x564925c93fd2 in function_leaky() /home/user/cmake-recipes/chapter-14/recipe-03/cxx-example/src/buggy.cpp:7
#2 0x564925c93fb2 in main /home/user/cmake-recipes/chapter-14/recipe-03/cxx-example/tests/leaky.cpp:4
#3 0x7ff98403df49 in __libc_start_main (/usr/lib/libc.so.6+0x20f49)
SUMMARY: AddressSanitizer: 8000 byte(s) leaked in 1 allocation(s).
\end{codebox}

相应地，我们可以直接运行\texttt{use\_after\_free}，得到详细的输出:

\begin{codebox}[shell]
$ ./build/tests/use_after_free

=================================================================
==18571==ERROR: AddressSanitizer: heap-use-after-free on address 0x6250000004d8 at pc 0x557ffa8b0102 bp 0x7ffe8c560200 sp 0x7ffe8c5601f0
READ of size 8 at 0x6250000004d8 thread T0
#0 0x557ffa8b0101 in function_use_after_free() /home/user/cmake-recipes/chapter-14/recipe-03/cxx-example/src/buggy.cpp:28
#1 0x557ffa8affb2 in main /home/user/cmake-recipes/chapter-14/recipe-03/cxx-example/tests/use_after_free.cpp:4
#2 0x7ff1d6088f49 in __libc_start_main (/usr/lib/libc.so.6+0x20f49)
#3 0x557ffa8afec9 in _start (/home/user/cmake-recipes/chapter-14/recipe-03/cxx-example/build/tests/use_after_free+0xec9)
0x6250000004d8 is located 984 bytes inside of 8000-byte region [0x625000000100,0x625000002040)
freed by thread T0 here:
#0 0x7ff1d6ded5a9 in operator delete[](void*) /build/gcc/src/gcc/libsanitizer/asan/asan_new_delete.cc:128
#1 0x557ffa8afffa in function_use_after_free() /home/user/cmake-recipes/chapter-14/recipe-03/cxx-example/src/buggy.cpp:24
#2 0x557ffa8affb2 in main /home/user/cmake-recipes/chapter-14/recipe-03/cxx-example/tests/use_after_free.cpp:4
#3 0x7ff1d6088f49 in __libc_start_main (/usr/lib/libc.so.6+0x20f49)
previously allocated by thread T0 here:
#0 0x7ff1d6dec669 in operator new[](unsigned long) /build/gcc/src/gcc/libsanitizer/asan/asan_new_delete.cc:82
#1 0x557ffa8affea in function_use_after_free() /home/user/cmake-recipes/chapter-14/recipe-03/cxx-example/src/buggy.cpp:19
#2 0x557ffa8affb2 in main /home/user/cmake-recipes/chapter-14/recipe-03/cxx-example/tests/use_after_free.cpp:4
#3 0x7ff1d6088f49 in __libc_start_main (/usr/lib/libc.so.6+0x20f49)
SUMMARY: AddressSanitizer: heap-use-after-free /home/user/cmake-recipes/chapter-14/recipe-03/cxx-example/src/buggy.cpp:28 in function_use_after_free()
Shadow bytes around the buggy address:
0x0c4a7fff8040: fd fd fd fd fd fd fd fd fd fd fd fd fd fd fd fd
0x0c4a7fff8050: fd fd fd fd fd fd fd fd fd fd fd fd fd fd fd fd
0x0c4a7fff8060: fd fd fd fd fd fd fd fd fd fd fd fd fd fd fd fd
0x0c4a7fff8070: fd fd fd fd fd fd fd fd fd fd fd fd fd fd fd fd
0x0c4a7fff8080: fd fd fd fd fd fd fd fd fd fd fd fd fd fd fd fd
=>0x0c4a7fff8090: fd fd fd fd fd fd fd fd fd fd fd[fd]fd fd fd fd
0x0c4a7fff80a0: fd fd fd fd fd fd fd fd fd fd fd fd fd fd fd fd
0x0c4a7fff80b0: fd fd fd fd fd fd fd fd fd fd fd fd fd fd fd fd
0x0c4a7fff80c0: fd fd fd fd fd fd fd fd fd fd fd fd fd fd fd fd
0x0c4a7fff80d0: fd fd fd fd fd fd fd fd fd fd fd fd fd fd fd fd
0x0c4a7fff80e0: fd fd fd fd fd fd fd fd fd fd fd fd fd fd fd fd
Shadow byte legend (one shadow byte represents 8 application bytes):
Addressable: 00
Partially addressable: 01 02 03 04 05 06 07
Heap left redzone: fa
Freed heap region: fd
Stack left redzone: f1
Stack mid redzone: f2
Stack right redzone: f3
Stack after return: f5
Stack use after scope: f8
Global redzone: f9
Global init order: f6
Poisoned by user: f7
Container overflow: fc
Array cookie: ac
Intra object redzone: bb
ASan internal: fe
Left alloca redzone: ca
Right alloca redzone: cb
==18571==ABORTING
\end{codebox}

如果我们在没有AddressSanitizer的情况下进行测试(默认情况下\texttt{ENABLE\_ASAN}是关闭的)，就不会报告错误：

\begin{codebox}[shell]
$ mkdir -p build_no_asan
$ cd build_no_asan
$ cmake ..
$ cmake --build .
$ cmake --build . --target test

Start 1: leaky
1/2 Test #1: leaky ............................ Passed 0.00 sec
Start 2: use_after_free
2/2 Test #2: use_after_free ................... Passed 0.00 sec
100% tests passed, 0 tests failed out of 2
\end{codebox}

实际上，泄漏只会浪费内存，而\texttt{use\_after\_free}可能会导致未定义行为。调试这些问题的一种方法是使用valgrind (http://valgrind.org )。

与前两个示例相反，我们使用了CTest脚本来配置、构建和测试代码，并将报告提交到面板。要了解此示例的工作原理，请仔细查看\texttt{ dashboard.cmake }脚本。首先，我们定义项目名称并设置主机报告和构建名称:

\begin{codebox}[cmake]
set(CTEST_PROJECT_NAME "example")
cmake_host_system_information(RESULT _site QUERY HOSTNAME)
set(CTEST_SITE ${_site})
set(CTEST_BUILD_NAME "${CMAKE_SYSTEM_NAME}-${CMAKE_HOST_SYSTEM_PROCESSOR}")
\end{codebox}

我们的例子中，\texttt{CTEST\_BUILD\_NAME}的计算结果是\texttt{Linux-x86\_64}。不同的操作系统下，可能会观察到不同的结果。

接下来，我们为源和构建目录指定路径:

\begin{codebox}[cmake]
set(CTEST_SOURCE_DIRECTORY "${CTEST_SCRIPT_DIRECTORY}")
set(CTEST_BINARY_DIRECTORY "${CTEST_SCRIPT_DIRECTORY}/build")
\end{codebox}

我们可以将生成器设置为\texttt{Unix Makefile}:

\begin{codebox}[cmake]
set(CTEST_CMAKE_GENERATOR "Unix Makefiles")
\end{codebox}

但是，对于更具可移植性的测试脚本，我们更愿意通过命令行提供生成器:

\begin{codebox}[shell]
$ ctest -S dashboard.cmake -D CTEST_CMAKE_GENERATOR="Unix Makefiles"
\end{codebox}

\texttt{ dashboard.cmake}中的下一个代码片段，将计算出机器上可用的CPU芯数量，并将测试步骤的并行级设置为可用CPU芯数量，以使总测试时间最小化:

\begin{codebox}[cmake]
include(ProcessorCount)
ProcessorCount(N)
if(NOT N EQUAL 0)
    set(CTEST_BUILD_FLAGS -j${N})
    set(ctest_test_args ${ctest_test_args} PARALLEL_LEVEL ${N})
endif()
\end{codebox}

接下来，我们开始测试步骤并配置代码，将\texttt{ENABLE\_ASAN}设置为\texttt{ON}:

\begin{codebox}[cmake]
ctest_start(Experimental)

ctest_configure(
  OPTIONS
    -DENABLE_ASAN:BOOL=ON
  )
\end{codebox}

\texttt{ dashboard.cmake }其他命令为映射到构建、测试、内存检查和提交步骤:

\begin{codebox}[cmake]
ctest_build()
ctest_test()

set(CTEST_MEMORYCHECK_TYPE "AddressSanitizer")

ctest_memcheck()
ctest_submit()
\end{codebox}

\mySmallsection{更多信息}

细心的读者会注意到，在链接目标之前，我们没有在系统上搜索AddressSanitizer。实际中，库查找工作已经提前做完，以避免在链接阶段出现意外。

有关AddressSanitizer文档和示例的更多信息，请参见\texthref{https://github.com/google/sanitizers/wiki/AddressSanitizer} 。AddressSanitizer并不仅限于C和C++。对于Fortran示例，读者可以参考 \texthref{https://github.com/dev-cafe/cmake-cookbook/tree/v1.0/chapter-14/recipe-03/fortran-example} 。

\begin{myInfoBox}
\textbf{NOTE}：可以在\texthref{https://github.com/arsenm/sanitizers-cmake} 上找到CMake程序，用来查找杀毒程序和调整编译器标志
\end{myInfoBox}

下面的博客文章讨论了如何添加对动态分析工具的支持，对我们很有启发性：\texthref{https://blog.kitware.com/ctest-cdash-add-support-for-new-dynamic-analysis-tools/}


% --- recipe 4: 14.4-chinese.md ---
\section{14.4 使用ThreadSaniiser向CDash报告数据争用}

\begin{myInfoBox}
\textbf{NOTE}：此示例代码可以在 \texthref{https://github.com/dev-cafe/cmake-cookbook/tree/v1.0/chapter-14/recipe-03} 中找到，其中包含一个C++示例。该示例在CMake 3.5版(或更高版本)中是有效的，并且已经在GNU/Linux、macOS和Windows上进行过测试。
\end{myInfoBox}

在这个示例中，我们将重用前一个示例中的方法，但是使用ThreadSanitizer或TSan，结合CTest和CDash，来检查数据竞争，并将它们报告给CDash。ThreadSanitizer的文档可以在网上找到，\texthref{https://github.com/google/sanitizers/wiki/ThreadSanitizerCppManual}

\mySmallsection{准备工作}

这个示例中，我们将使用以下示例代码(\texttt{example.cpp}):

\begin{codebox}[cpp]
#include <chrono>
#include <iostream>
#include <thread>

static const int num_threads = 16;

void increase(int i, int &s) {
  std::this_thread::sleep_for(std::chrono::seconds(1));
  std::cout << "thread " << i << " increases " << s++ << std::endl;
}

int main() {
  std::thread t[num_threads];
  
  int s = 0;
  
  // start threads
  for (auto i = 0; i < num_threads; i++) {
    t[i] = std::thread(increase, i, std::ref(s));
  }
  
  // join threads with main thread
  for (auto i = 0; i < num_threads; i++) {
    t[i].join();
  }
  
  std::cout << "final s: " << s << std::endl;
  
  return 0;
}
\end{codebox}

这个示例代码中，我们启动16个线程，每个线程都调用\texttt{increase}函数。\texttt{increase}函数休眠1s，然后打印并递增一个整数\texttt{s}。我们预计此示例代码将显示数据竞争，因为所有线程读取和修改相同的地址，而不需要任何显式同步或协调。换句话说，我们期望在代码末尾打印的最终\texttt{s}，每次的结果都不同。代码有bug，我们将尝试在ThreadSanitizer的帮助下识别数据竞争。如果不运行ThreadSanitizer，我们可能不会看到代码有任何问题:

\begin{codebox}[shell]
$ ./example

thread thread 0 increases 01 increases 1
thread 9 increases 2
thread 4 increases 3
thread 10 increases 4
thread 2 increases 5
thread 3 increases 6
thread 13 increases 7
thread thread 7 increases 8
thread 14 increases 9
thread 8 increases 10
thread 12 increases 11
thread 15 increases 12
thread 11 increases 13

5 increases 14
thread 6 increases 15
final s: 16
\end{codebox}

\mySmallsection{具体实施}

\begin{enumerate}
\item 文件\texttt{CMakeLists.txt}首先定义一个受支持的最低版本、项目名称、受支持的语言。在本例中，定义了C++11标准项目:

\begin{codebox}[cmake]
cmake_minimum_required(VERSION 3.5 FATAL_ERROR)
project(recipe-04 LANGUAGES CXX)
set(CMAKE_CXX_STANDARD 11)
set(CMAKE_CXX_EXTENSIONS OFF)
set(CMAKE_CXX_STANDARD_REQUIRED ON)
\end{codebox}

\item 接下来，找到线程库，定义可执行文件，并将其链接到线程库:

\begin{codebox}[cmake]
find_package(Threads REQUIRED)
add_executable(example example.cpp)
target_link_libraries(example
  PUBLIC
    Threads::Threads
  )
\end{codebox}

\item 然后，提供编译选项和代码，并链接到ThreadSanitizer:

\begin{codebox}[cmake]
option(ENABLE_TSAN "Enable ThreadSanitizer" OFF)
if(ENABLE_TSAN)
  if(CMAKE_CXX_COMPILER_ID MATCHES GNU)
    message(STATUS "ThreadSanitizer enabled")
    target_compile_options(example
    PUBLIC
        -g -O1 -fsanitize=thread -fno-omit-frame-pointer -fPIC
    )
    target_link_libraries(example
      PUBLIC
        tsan
      )
  else()
    message(WARNING "ThreadSanitizer not supported for this compiler")
  endif()
endif()
\end{codebox}

\item 最后，编译测试用例:

\begin{codebox}[cmake]
enable_testing()
# allow to report to a cdash dashboard
include(CTest)
add_test(
  NAME
    example
  COMMAND
    $<TARGET_FILE:example>
  )
\end{codebox}

\item \texttt{ CTestConfig.cmake}没有变化:

\begin{codebox}[cmake]
set(CTEST_DROP_METHOD "http")
set(CTEST_DROP_SITE "my.cdash.org")
set(CTEST_DROP_LOCATION "/submit.php?project=cmake-cookbook")
set(CTEST_DROP_SITE_CDASH TRUE)
\end{codebox}

\item \texttt{dashboard.cmake }需要为TSan进行简单修改:

\begin{codebox}[cmake]
set(CTEST_PROJECT_NAME "example")
cmake_host_system_information(RESULT _site QUERY HOSTNAME)
set(CTEST_SITE ${_site})
set(CTEST_BUILD_NAME "${CMAKE_SYSTEM_NAME}-${CMAKE_HOST_SYSTEM_PROCESSOR}")
set(CTEST_SOURCE_DIRECTORY "${CTEST_SCRIPT_DIRECTORY}")
set(CTEST_BINARY_DIRECTORY "${CTEST_SCRIPT_DIRECTORY}/build")
include(ProcessorCount)
ProcessorCount(N)
if(NOT N EQUAL 0)
  set(CTEST_BUILD_FLAGS -j${N})
  set(ctest_test_args ${ctest_test_args} PARALLEL_LEVEL ${N})
endif()
ctest_start(Experimental)
ctest_configure(
  OPTIONS
    -DENABLE_TSAN:BOOL=ON
  )
ctest_build()
ctest_test()
set(CTEST_MEMORYCHECK_TYPE "ThreadSanitizer")
ctest_memcheck()
ctest_submit()
\end{codebox}

\item 让我们以这个例子为例。通过\texttt{CTEST\_CMAKE\_GENERATOR}选项来设置生成器:

\begin{codebox}[shell]
$ ctest -S dashboard.cmake -D CTEST_CMAKE_GENERATOR="Unix Makefiles"
Each . represents 1024 bytes of output
. Size of output: 0K
Each symbol represents 1024 bytes of output.
'!' represents an error and '*' a warning.
. Size of output: 0K
\end{codebox}

\item 在面板上，我们将看到以下内容:

\myGraphic{0.9}{images/chapter14/14-8.png}{}

\item 我们可以看到更详细的动态分析:

\myGraphic{0.9}{images/chapter14/14-9.png}{}

\end{enumerate}

\mySmallsection{工作原理}

该示例\texttt{CMakeLists.txt}的核心部分:

\begin{codebox}[cmake]
option(ENABLE_TSAN "Enable ThreadSanitizer" OFF)

if(ENABLE_TSAN)
  if(CMAKE_CXX_COMPILER_ID MATCHES GNU)
    message(STATUS "ThreadSanitizer enabled")
    target_compile_options(example
    PUBLIC
        -g -O1 -fsanitize=thread -fno-omit-frame-pointer -fPIC
    )
    target_link_libraries(example
      PUBLIC
        tsan
      )
  else()
    message(WARNING "ThreadSanitizer not supported for this compiler")
  endif()
endif()
\end{codebox}

\texttt{dashboard.cmake}也需要更新:

\begin{codebox}[cmake]
# ...

ctest_start(Experimental)

ctest_configure(
  OPTIONS
    -DENABLE_TSAN:BOOL=ON
  )

ctest_build()
ctest_test()

set(CTEST_MEMORYCHECK_TYPE "ThreadSanitizer")
ctest_memcheck()

ctest_submit()
\end{codebox}

和上一个示例一样，我们也可以在本地查看ThreadSanitizer的输出:

\begin{codebox}[shell]
$ mkdir -p build
$ cd build
$ cmake -DENABLE_TSAN=ON ..
$ cmake --build .
$ cmake --build . --target test

Start 1: example
1/1 Test #1: example ..........................***Failed 1.07 sec
0% tests passed, 1 tests failed out of 1
$ ./build/example
thread 0 increases 0
==================
WARNING: ThreadSanitizer: data race (pid=24563)
... lots of output ...
SUMMARY: ThreadSanitizer: data race /home/user/cmake-recipes/chapter-14/recipe-04/cxx-example/example
\end{codebox}

\mySmallsection{更多信息}

对使用OpenMP的应用TSan是很常见的，但是请注意，在某些情况下，OpenMP会在TSan下生成误检的结果。对于Clang编译器，一个解决方案是用\texttt{-DLIBOMP\_TSAN\_SUPPORT=TRUE}重新编译编译器本身及其\texttt{libomp}。通常，以合理的方式使用TSan可能需要重新编译整个工具堆栈，以避免误报。在使用pybind11的C++项目的情况，我们可能需要重新编译Python，并启用TSan来获得有意义的东西。或者，Python绑定可以通过使用TSan抑制而被排除在外，如 \texthref{https://github.com/google/sanitizers/wiki/threadsanitizersuppression} 。例如：如果一个动态库同时被一个经过TSan的二进制文件和一个Python插件调用，那么这种情况可能是不可能使用TSan。

下面的博客文章讨论了如何添加对动态分析工具的支持：\texthref{https://blog.kitware.com/ctest-cdash-add-support-for-new-dynamic-analysis-tools/}


